% ============================================================================
% CHAPITRE 8 : FORMULAIRE SYNTHÉTIQUE ET AIDE-MÉMOIRE (VERSION IMPRIMABLE 4 PAGES)
% ============================================================================
\cleardoublepage
\newgeometry{top=1.3cm,bottom=1.4cm,left=1.4cm,right=1.4cm}
\refstepcounter{chapter}
\addcontentsline{toc}{chapter}{\protect\numberline{\thechapter}Formulaire Synthétique et Aide-Mémoire (Imprimable -- 4 pages)}
\chaptermark{Formulaire Synthétique et Aide-Mémoire}

% ============================================================================
% PAGE 1/4 : FORMULAIRE DE TRIGONOMÉTRIE (1 PAGE A4 IMPRIMABLE)
% ============================================================================
\thispagestyle{plain}
\refstepcounter{section}
\addcontentsline{toc}{section}{\protect\numberline{\thesection}Formulaire de Trigonométrie}

\begin{center}
    {\LARGE\bfseries\color{uclblue} Formulaire de Trigonométrie}\\[2pt]
    {\small\color{darkgray} Examen d'entrée Ingénieur Civil -- UCLouvain (Fiche récapitulative imprimable 1/4)}
    \vspace{2pt}
    \hrule height 1pt
\end{center}
\vspace{-6pt}

{\fontsize{9.5pt}{11.5pt}\selectfont
\setlength{\columnsep}{18pt}
\begin{multicols}{2}

\noindent{\large\bfseries\color{uclblue} 1. Identités Fondamentales}
\begin{itemize}[leftmargin=1.2em, itemsep=1.5pt, topsep=2pt]
    \item $\cos^2\alpha + \sin^2\alpha = 1$
    \item $\tan\alpha = \dfrac{\sin\alpha}{\cos\alpha} \;\left(\alpha \neq \frac{\pi}{2}+k\pi\right)$, \quad $\cot\alpha = \dfrac{\cos\alpha}{\sin\alpha} \;(\alpha \neq k\pi)$
    \item $1 + \tan^2\alpha = \dfrac{1}{\cos^2\alpha}$ \quad et \quad $1 + \cot^2\alpha = \dfrac{1}{\sin^2\alpha}$
\end{itemize}

\vspace{3pt}
\noindent{\large\bfseries\color{uclblue} 2. Angles Associés}
\begin{itemize}[leftmargin=1.2em, itemsep=1.5pt, topsep=2pt]
    \item \textbf{Opposés ($-\alpha$) :} $\cos(-\alpha) = \cos\alpha$, \; $\sin(-\alpha) = -\sin\alpha$, \; $\tan(-\alpha) = -\tan\alpha$
    \item \textbf{Supplémentaires ($\pi-\alpha$) :} $\cos(\pi-\alpha) = -\cos\alpha$, \; $\sin(\pi-\alpha) = \sin\alpha$, \; $\tan(\pi-\alpha) = -\tan\alpha$
    \item \textbf{Antisupplémentaires ($\pi+\alpha$) :} $\cos(\pi+\alpha) = -\cos\alpha$, \; $\sin(\pi+\alpha) = -\sin\alpha$, \; $\tan(\pi+\alpha) = \tan\alpha$
    \item \textbf{Complémentaires ($\frac{\pi}{2}-\alpha$) :} $\cos\left(\frac{\pi}{2}-\alpha\right) = \sin\alpha$, \; $\sin\left(\frac{\pi}{2}-\alpha\right) = \cos\alpha$, \; $\tan\left(\frac{\pi}{2}-\alpha\right) = \cot\alpha$
    \item \textbf{Anticomplémentaires ($\frac{\pi}{2}+\alpha$) :} $\cos\left(\frac{\pi}{2}+\alpha\right) = -\sin\alpha$, \; $\sin\left(\frac{\pi}{2}+\alpha\right) = \cos\alpha$, \; $\tan\left(\frac{\pi}{2}+\alpha\right) = -\cot\alpha$
\end{itemize}

\vspace{3pt}
\noindent{\large\bfseries\color{uclblue} 3. Formules d'Addition}
\begin{itemize}[leftmargin=1.2em, itemsep=1.5pt, topsep=2pt]
    \item $\cos(\alpha \pm \beta) = \cos\alpha\cos\beta \mp \sin\alpha\sin\beta$
    \item $\sin(\alpha \pm \beta) = \sin\alpha\cos\beta \pm \cos\alpha\sin\beta$
    \item $\tan(\alpha \pm \beta) = \dfrac{\tan\alpha \pm \tan\beta}{1 \mp \tan\alpha\tan\beta}$
\end{itemize}

\vspace{3pt}
\noindent{\large\bfseries\color{uclblue} 4. Formules de Duplication, Carnot et Angle Demi}
\begin{itemize}[leftmargin=1.2em, itemsep=1.5pt, topsep=2pt]
    \item $\cos(2\alpha) = \cos^2\alpha - \sin^2\alpha = 2\cos^2\alpha - 1 = 1 - 2\sin^2\alpha$
    \item $\sin(2\alpha) = 2\sin\alpha\cos\alpha$ \quad et \quad $\tan(2\alpha) = \dfrac{2\tan\alpha}{1 - \tan^2\alpha}$
    \item \textbf{Carnot :} $\cos^2\alpha = \dfrac{1+\cos(2\alpha)}{2}$, \quad $\sin^2\alpha = \dfrac{1-\cos(2\alpha)}{2}$
    \item \textbf{Arc moitié ($t = \tan(\alpha/2)$) :}\\
    $\cos\alpha = \dfrac{1-t^2}{1+t^2}, \quad \sin\alpha = \dfrac{2t}{1+t^2}, \quad \tan\alpha = \dfrac{2t}{1-t^2}$
\end{itemize}

\vspace{3pt}
\noindent{\large\bfseries\color{uclblue} 5. Formules de Simpson (Factorisation)}
\begin{itemize}[leftmargin=1.2em, itemsep=2pt, topsep=2pt]
    \item $\cos p + \cos q = 2\cos\left(\dfrac{p+q}{2}\right)\cos\left(\dfrac{p-q}{2}\right)$
    \item $\cos p - \cos q = -2\sin\left(\dfrac{p+q}{2}\right)\sin\left(\dfrac{p-q}{2}\right)$
    \item $\sin p + \sin q = 2\sin\left(\dfrac{p+q}{2}\right)\cos\left(\dfrac{p-q}{2}\right)$
    \item $\sin p - \sin q = 2\sin\left(\dfrac{p-q}{2}\right)\cos\left(\dfrac{p+q}{2}\right)$
    \item $\tan p \pm \tan q = \dfrac{\sin(p \pm q)}{\cos p \cos q}$
\end{itemize}

\columnbreak

\noindent{\large\bfseries\color{uclblue} 6. Arcs, Secteurs et Triangles}

\vspace{2pt}
\noindent\textbf{Cercle de rayon $R$ ($\theta$ en radians) :}
\begin{itemize}[leftmargin=1.2em, itemsep=1.5pt, topsep=2pt]
    \item Longueur d'un arc : $L = R\theta$ \quad (Corde : $c = 2R\sin(\theta/2)$)
    \item Aire d'un secteur : $A_{\text{secteur}} = \dfrac{1}{2}R^2\theta$
    \item Aire d'un segment : $A_{\text{segment}} = \dfrac{1}{2}R^2(\theta - \sin\theta)$
\end{itemize}

\vspace{2pt}
\noindent\textbf{Triangle Rectangle (SOH-CAH-TOA) :}
\begin{itemize}[leftmargin=1.2em, itemsep=1.5pt, topsep=2pt]
    \item $\sin\theta = \dfrac{\text{Opposé}}{\text{Hypoténuse}}$
    \item $\cos\theta = \dfrac{\text{Adjacent}}{\text{Hypoténuse}}$
    \item $\tan\theta = \dfrac{\text{Opposé}}{\text{Adjacent}}$
\end{itemize}

\vspace{2pt}
\noindent\textbf{Triangle Quelconque ($A+B+C = \pi$) :}
\begin{itemize}[leftmargin=1.2em, itemsep=2pt, topsep=2pt]
    \item \textbf{Loi des sinus :} $\dfrac{a}{\sin A} = \dfrac{b}{\sin B} = \dfrac{c}{\sin C} = 2R$
    \item \textbf{Loi des cosinus (Al-Kashi) :}\\
    $a^2 = b^2 + c^2 - 2bc\cos A$
    \item \textbf{Aire :} $S = \dfrac{1}{2}bc\sin A = p\cdot r = \sqrt{p(p-a)(p-b)(p-c)}$\\
    avec $p = \dfrac{a+b+c}{2}$ (demi-périmètre) et $r$ rayon inscrit.
\end{itemize}

\vspace{4pt}
\noindent{\large\bfseries\color{uclblue} 7. Astuce d'Euler (Complexes)}
La formule d'Euler $e^{i\alpha} = \cos\alpha + i\sin\alpha$ permet de retrouver les formules d'addition, de duplication et de triplication en développant :
\[
    e^{i(\alpha+\beta)} = e^{i\alpha}e^{i\beta} \qquad \text{et} \qquad e^{in\alpha} = (e^{i\alpha})^n
\]

\vspace{2pt}
\noindent{\large\bfseries\color{uclblue} 8. Angles Remarquables}
\begin{center}
\vspace{-2pt}
\renewcommand{\arraystretch}{1.45}
\setlength{\tabcolsep}{4.5pt}
\begin{tabular}{|c|c|c|c|c|c|c|}
\hline
$\theta$ & $0$ & $\frac{\pi}{12}$ & $\frac{\pi}{6}$ & $\frac{\pi}{4}$ & $\frac{\pi}{3}$ & $\frac{\pi}{2}$ \\ \hline
$\sin$ & $0$ & $\frac{\sqrt{6}-\sqrt{2}}{4}$ & $\frac{1}{2}$ & $\frac{\sqrt{2}}{2}$ & $\frac{\sqrt{3}}{2}$ & $1$ \\ \hline
$\cos$ & $1$ & $\frac{\sqrt{6}+\sqrt{2}}{4}$ & $\frac{\sqrt{3}}{2}$ & $\frac{\sqrt{2}}{2}$ & $\frac{1}{2}$ & $0$ \\ \hline
$\tan$ & $0$ & $2-\sqrt{3}$ & $\frac{\sqrt{3}}{3}$ & $1$ & $\sqrt{3}$ & $-$ \\ \hline
\end{tabular}
\end{center}
\vspace{-2pt}
\noindent\emph{Note :} Pour $\theta = \frac{\pi}{8}$ ($22,5\degree$), $\tan\left(\frac{\pi}{8}\right) = \sqrt{2}-1$.

\end{multicols}
}

\newpage
% ============================================================================
% PAGE 2/4 : ÉQUATIONS TRIGONOMÉTRIQUES : TECHNIQUES ET EXEMPLES (1 PAGE A4)
% ============================================================================
\thispagestyle{plain}
\refstepcounter{section}
\addcontentsline{toc}{section}{\protect\numberline{\thesection}Équations Trigonométriques : Techniques et Exemples}

\begin{center}
    {\LARGE\bfseries\color{uclblue} Équations Trigonométriques : Techniques et Exemples}\\[2pt]
    {\small\color{darkgray} Examen d'entrée Ingénieur Civil -- UCLouvain (Fiche récapitulative imprimable 2/4)}
    \vspace{2pt}
    \hrule height 1pt
\end{center}
\vspace{-6pt}

{\fontsize{9.5pt}{11.5pt}\selectfont
\setlength{\columnsep}{18pt}
\begin{multicols}{2}

\noindent{\large\bfseries\color{uclblue} 1. Équations Élémentaires}

\vspace{2pt}
\noindent Règles de base ($k \in \mathbb{Z}$) :
\begin{itemize}[leftmargin=1.2em, itemsep=2pt, topsep=2pt]
    \item $\sin(\alpha) = \sin(\beta) \iff$\\
    $\alpha = \beta + 2k\pi \quad\text{ou}\quad \alpha = \pi - \beta + 2k\pi$
    \item $\cos(\alpha) = \cos(\beta) \iff$\\
    $\alpha = \pm\beta + 2k\pi$
    \item $\tan(\alpha) = \tan(\beta) \iff$\\
    $\alpha = \beta + k\pi \quad \left(\text{CE : } \alpha,\beta \neq \frac{\pi}{2}+m\pi\right)$
\end{itemize}
\textit{Astuce :} Pour $\sin(\alpha) = \cos(\beta)$, on transforme en $\sin(\alpha) = \sin\left(\frac{\pi}{2} - \beta\right)$ ou $\cos\left(\frac{\pi}{2} - \alpha\right) = \cos(\beta)$.

\vspace{6pt}
\noindent{\large\bfseries\color{uclblue} 2. Polynômes et Substitution}

\vspace{2pt}
\noindent Pour des équations avec des puissances paires, on peut remplacer et poser une variable.\\
\textbf{Exemple (UCLouvain, Juillet 2025) :}\\
$\sin^6 x - \cos^6 x = 3\sin^2 x \cos^2 x$
\begin{itemize}[leftmargin=1.2em, itemsep=2pt, topsep=2pt]
    \item Poser $y = \sin^2 x$, ce qui donne $\cos^2 x = 1 - y$.
    \item Remplacer : $y^3 - (1-y)^3 = 3y(1-y)$.
    \item Simplifier pour obtenir une équation algébrique en $y$ :
    \[
        2y^3 - 1 = 0 \implies \sin^2 x = \sqrt[3]{\frac{1}{2}}
    \]
\end{itemize}
\textit{Astuce :} Utiliser l'identité $a^3 \mp b^3 = (a\mp b)(a^2 \pm ab + b^2)$ pour factoriser quand c'est pertinent, ou les identités classiques $\sin^4 x + \cos^4 x = 1 - 2\sin^2 x\cos^2 x$ et $\sin^6 x + \cos^6 x = 1 - 3\sin^2 x\cos^2 x$.

\vspace{6pt}
\noindent{\large\bfseries\color{uclblue} 3. Linéaire en $\sin x, \cos x$}

\vspace{2pt}
\noindent Forme : $a\sin x + b\cos x = c \quad (\text{solutions ssi } c^2 \le a^2+b^2)$
\begin{itemize}[leftmargin=1.2em, itemsep=2pt, topsep=2pt]
    \item Diviser par $a$ : $\sin x + \dfrac{b}{a}\cos x = \dfrac{c}{a}$.
    \item Poser $\tan(y) = \dfrac{b}{a}$. L'équation devient :
    \[
        \sin(x)\cos(y) + \sin(y)\cos(x) = \frac{c}{a}\cos(y)
    \]
    \item Ce qui équivaut, avec la formule d'addition, à :
    \[
        \sin(x+y) = \frac{c}{a}\cos(y)
    \]
\end{itemize}

\columnbreak

\noindent{\large\bfseries\color{uclblue} 4. Homogène en $\sin x, \cos x$}

\vspace{2pt}
\noindent Polynôme en $\sin, \cos$ où la somme des exposants de chaque terme est constante.\\
\textbf{Exemple (UCLouvain, Juillet 2018) :}\\
$\sin^3 x + 2\cos^3 x = 3\sin^2 x \cos x$ (Degré 3)
\begin{itemize}[leftmargin=1.2em, itemsep=2pt, topsep=2pt]
    \item Vérifier qu'on ne peut pas mettre de terme en évidence.
    \item Vérifier si $\cos x = 0$ est solution, puis diviser les deux membres par $\cos^k x$ (ici $k=3$) pour se ramener à une équation en $\tan x$.
    \item On obtient l'équation algébrique :
    \[
        \tan^3 x + 2 = 3\tan^2 x
    \]
\end{itemize}

\vspace{6pt}
\noindent{\large\bfseries\color{uclblue} 5. Équation Symétrique}

\vspace{2pt}
\noindent L'équation ne change pas si on inverse $\sin x$ et $\cos x$.\\
\textbf{Exemple (UCLouvain, Juillet 2022) :}\\
$\cos x + \sin x = \sqrt{2}\cos x \sin x$
\begin{itemize}[leftmargin=1.2em, itemsep=2pt, topsep=2pt]
    \item Technique clé : poser $\mathbf{x = y + \dfrac{\pi}{4}}$.
    \item On sait via les formules d'addition que :
    \[
        \sin x = \frac{\sqrt{2}}{2}(\cos y + \sin y), \quad \cos x = \frac{\sqrt{2}}{2}(\cos y - \sin y)
    \]
    \item Remplacer dans l'équation ($\sin x + \cos x = \sqrt{2}\cos y$ et $\sin x\cos x = \cos^2 y - \frac{1}{2}$) pour obtenir une équation du second degré en $\cos y$ facile à résoudre.
\end{itemize}

\vspace{6pt}
\noindent{\large\bfseries\color{uclblue} 6. Formules de Simpson}

\vspace{2pt}
\noindent Pour les sommes/différences de sinus et cosinus avec des multiples ($2x, 3x\dots$).\\
\textbf{Exemple (UCLouvain, Septembre 2019) :}\\
$\sin(x) - 2\sin(3x) + \sin(5x) = -\cos(2x) + \cos(4x)$
\begin{itemize}[leftmargin=1.2em, itemsep=2pt, topsep=2pt]
    \item Regrouper stratégiquement les termes :
    \[
        \sin(5x) + \sin(x) = 2\sin(3x)\cos(2x)
    \]
    \item Regrouper :
    \[
        \cos(4x) - \cos(2x) = -2\sin(3x)\sin(x)
    \]
    \item L'équation devient factorisable par $\sin(3x)$ :
    \[
        2\sin(3x)\bigl(\cos(2x) - 1 + \sin(x)\bigr) = 0
    \]
    donc soit $\sin(3x) = 0$, soit l'autre facteur est nul.
\end{itemize}

\end{multicols}
}

\newpage
% ============================================================================
% PAGE 3/4 : GÉOMÉTRIE SYNTHÉTIQUE & ANALYTIQUE PLANE (CONIQUES) (1 PAGE A4)
% ============================================================================
\thispagestyle{plain}
\refstepcounter{section}
\addcontentsline{toc}{section}{\protect\numberline{\thesection}Formulaire de Géométrie Synthétique et Analytique Plane (Coniques)}

\begin{center}
    {\LARGE\bfseries\color{uclblue} Géométrie Synthétique \& Analytique Plane (Coniques)}\\[2pt]
    {\small\color{darkgray} Examen d'entrée Ingénieur Civil -- UCLouvain (Fiche récapitulative imprimable 3/4)}
    \vspace{2pt}
    \hrule height 1pt
\end{center}
\vspace{-6pt}

{\fontsize{9.2pt}{11.2pt}\selectfont
\setlength{\columnsep}{18pt}
\begin{multicols}{2}

\noindent{\large\bfseries\color{uclblue} 1. Géométrie Synthétique du Cercle}
\begin{itemize}[leftmargin=1.2em, itemsep=1.5pt, topsep=2pt]
    \item \textbf{Médiatrice d'une corde $[AB]$ :} La droite reliant le centre $O$ au milieu $M$ de $[AB]$ est perpendiculaire à $[AB]$ ($|OM|^2 + |AM|^2 = r^2$).
    \item \textbf{Angle inscrit et angle au centre} interceptant $\wideparen{AB}$ :
    \[
        |\widehat{AMB}| = \frac{1}{2}|\widehat{AOB}|
    \]
    \item \textbf{Demi-cercle :} Si $[AB]$ est un diamètre et $M \in \mathcal{C}$, alors $\widehat{AMB} = 90\degree$.
    \item \textbf{Quadrilatère inscriptible $ABCD$ :} Les angles opposés sont supplémentaires ($|\widehat{A}| + |\widehat{C}| = 180\degree = |\widehat{B}| + |\widehat{D}|$).
    \item \textbf{Puissance d'un point $P$ par rapport à $\mathcal{C}(O,r)$ :}
    \begin{itemize}[leftmargin=1em, itemsep=1pt, topsep=1pt]
        \item $P$ extérieur : $|PA|\cdot|PB| = |PT|^2 = |PO|^2 - r^2$
        \item $P$ intérieur : $|PA|\cdot|PB| = r^2 - |PO|^2$
    \end{itemize}
\end{itemize}

\vspace{3pt}
\noindent{\large\bfseries\color{uclblue} 2. Relations Métriques \& Aires}
\begin{itemize}[leftmargin=1.2em, itemsep=1.5pt, topsep=2pt]
    \item \textbf{Triangle rectangle en $C$} (hauteur $h = |CH|$ sur l'hypoténuse $c = p + q$) : $h^2 = p\cdot q$, \; $b^2 = c\cdot p$, \; $a^2 = c\cdot q$.
    \item \textbf{Figures semblables} de rapport $\rho = \frac{L_2}{L_1}$ : $\dfrac{\text{Aire}_2}{\text{Aire}_1} = \rho^2$.
    \item \textbf{Triangles de même hauteur :} $\dfrac{\text{Aire}_1}{\text{Aire}_2} = \dfrac{\text{base}_1}{\text{base}_2}$.
    \item \textbf{Identité du parallélogramme :}\\
    $|AC|^2 + |BD|^2 = 2\bigl(|AB|^2 + |BC|^2\bigr)$.
    \item \textbf{Solides :} $V_{\text{cyl}} = \pi R^2 h$, \; $V_{\text{cône}} = \frac{1}{3}\pi R^2 h$, \; $V_{\text{pyr}} = \frac{1}{3}\mathcal{B}h$, \; $V_{\text{sphère}} = \frac{4}{3}\pi r^3$, \; $A_{\text{sphère}} = 4\pi r^2$.
\end{itemize}

\vspace{3pt}
\noindent{\large\bfseries\color{uclblue} 3. Vecteurs, Droites et Cercles dans $\R^2$}
\begin{itemize}[leftmargin=1.2em, itemsep=1.5pt, topsep=2pt]
    \item \textbf{Produit scalaire :} $\vec{u}\cdot\vec{v} = u_x v_x + u_y v_y = \|\vec{u}\|\|\vec{v}\|\cos\theta$.
    \item \textbf{Orthogonalité :} $\vec{u} \perp \vec{v} \iff u_x v_x + u_y v_y = 0$.\\
    Vecteur orthogonal à $(a,b)$ : $(-b, a)$.
    \item \textbf{Droite $d \equiv ax + by + c = 0$ :} normale $\vec{n} = (a,b)$, vecteur directeur $\vec{d} = (-b,a)$, pente $m = -\frac{a}{b}$.\\
    Perpendicularité de $y = m_1 x + p_1$ et $y = m_2 x + p_2$ : $m_1 m_2 = -1$.
    \item \textbf{Distance d'un point $P(x_0,y_0)$ à $d \equiv ax+by+c=0$ :}
    \[
        d(P, d) = \frac{|ax_0 + by_0 + c|}{\sqrt{a^2 + b^2}}
    \]
    \item \textbf{Cercle de centre $\Omega(x_0,y_0)$ et rayon $R$ :}\\
    $(x - x_0)^2 + (y - y_0)^2 = R^2$.\\
    Le centre du cercle circonscrit à $ABC$ est l'intersection des médiatrices de $[AB]$ et $[AC]$.
\end{itemize}

\columnbreak

\noindent{\large\bfseries\color{uclblue} 4. Les Coniques dans le Plan}

\vspace{2pt}
\noindent\textbf{Définition par foyer $F$, directrice $d$ et excentricité $e$ :}
\[
    \frac{d(M, F)}{d(M, d)} = e \quad \begin{cases} 0 < e < 1 & \text{: Ellipse} \\ e = 1 & \text{: Parabole} \\ e > 1 & \text{: Hyperbole} \end{cases}
\]

\vspace{2pt}
\noindent\textbf{A. Ellipse de centre $C(x_0, y_0)$ :}
\[
    \frac{(x-x_0)^2}{a^2} + \frac{(y-y_0)^2}{b^2} = 1
\]
\begin{itemize}[leftmargin=1.2em, itemsep=1.5pt, topsep=2pt]
    \item \textbf{Si $a > b$ (grand axe horizontal $2a$) :}\\
    $c = \sqrt{a^2 - b^2}$, \; foyers $F_{1,2}(x_0 \pm c,\; y_0)$, \; $e = \dfrac{c}{a} < 1$,\\
    directrices $x = x_0 \pm \dfrac{a^2}{c}$. \quad ($|MF_1| + |MF_2| = 2a$)
    \item \textbf{Si $b > a$ (grand axe vertical $2b$) :}\\
    $c = \sqrt{b^2 - a^2}$, \; foyers $F_{1,2}(x_0,\; y_0 \pm c)$, \; $e = \dfrac{c}{b} < 1$,\\
    directrices $y = y_0 \pm \dfrac{b^2}{c}$.
\end{itemize}

\vspace{3pt}
\noindent\textbf{B. Hyperbole de centre $C(x_0, y_0)$ (axe horizontal) :}
\[
    \frac{(x-x_0)^2}{a^2} - \frac{(y-y_0)^2}{b^2} = 1
\]
\begin{itemize}[leftmargin=1.2em, itemsep=1.5pt, topsep=2pt]
    \item Sommets : $S_{1,2}(x_0 \pm a,\; y_0)$.
    \item Distance focale : $\mathbf{c = \sqrt{a^2 + b^2}}$, \; foyers $F_{1,2}(x_0 \pm c,\; y_0)$.
    \item Excentricité : $e = \dfrac{c}{a} > 1$, \; directrices : $x = x_0 \pm \dfrac{a^2}{c}$.
    \item \textbf{Asymptotes :} $y - y_0 = \pm \dfrac{b}{a}(x - x_0)$.
\end{itemize}

\vspace{3pt}
\noindent\textbf{C. Parabole de sommet $S(x_0, y_0)$ :}
\begin{itemize}[leftmargin=1.2em, itemsep=1.5pt, topsep=2pt]
    \item \textbf{Axe vertical :} $(x-x_0)^2 = 2p(y-y_0) \iff y - y_0 = a(x-x_0)^2$ (avec $p = \frac{1}{2a}$).\\
    Foyer $F\left(x_0,\; y_0 + \frac{p}{2}\right) = \left(x_0,\; y_0 + \frac{1}{4a}\right)$,\\
    Directrice $d \equiv y = y_0 - \frac{p}{2} = y_0 - \frac{1}{4a}$.
    \item \textbf{Axe horizontal :} $(y-y_0)^2 = 2p(x-x_0)$, foyer $F\left(x_0 + \frac{p}{2}, y_0\right)$, directrice $x = x_0 - \frac{p}{2}$.
\end{itemize}

\vspace{3pt}
\noindent{\large\bfseries\color{uclblue} 5. Tangentes à une Conique}
\begin{itemize}[leftmargin=1.2em, itemsep=1.5pt, topsep=2pt]
    \item \textbf{Condition de tangence d'une droite $y = mx+p$ :} Substituer $y$ dans l'équation de la conique et imposer que le discriminant de l'équation du $2^{\text{e}}$ degré en $x$ soit nul : $\mathbf{\Delta = B^2 - 4AC = 0}$.
    \item \textbf{Vecteur tangent à $y = f(x)$ en $x_1$ :} $\vec{t} = \bigl(1,\; f'(x_1)\bigr)$.
\end{itemize}

\end{multicols}
}

\newpage
% ============================================================================
% PAGE 4/4 : GÉOMÉTRIE ANALYTIQUE ET VECTORIELLE DANS L'ESPACE (1 PAGE A4)
% ============================================================================
\thispagestyle{plain}
\refstepcounter{section}
\addcontentsline{toc}{section}{\protect\numberline{\thesection}Formulaire de Géométrie Analytique et Vectorielle dans l'Espace}

\begin{center}
    {\LARGE\bfseries\color{uclblue} Géométrie Analytique et Vectorielle dans l'Espace ($\R^3$)}\\[2pt]
    {\small\color{darkgray} Examen d'entrée Ingénieur Civil -- UCLouvain (Fiche récapitulative imprimable 4/4)}
    \vspace{2pt}
    \hrule height 1pt
\end{center}
\vspace{-6pt}

{\fontsize{9.5pt}{11.5pt}\selectfont
\setlength{\columnsep}{18pt}
\begin{multicols}{2}

\noindent{\large\bfseries\color{uclblue} 1. Outils Vectoriels dans $\R^3$}
\begin{itemize}[leftmargin=1.2em, itemsep=2pt, topsep=2pt]
    \item \textbf{Distance et milieu de $A(x_A,y_A,z_A)$ et $B(x_B,y_B,z_B)$ :}
    \[
        |AB| = \sqrt{(x_B-x_A)^2 + (y_B-y_A)^2 + (z_B-z_A)^2}
    \]
    \[
        M = \left(\frac{x_A+x_B}{2},\; \frac{y_A+y_B}{2},\; \frac{z_A+z_B}{2}\right)
    \]
    \item \textbf{Produit scalaire et angle $\theta$ :}
    \[
        \vec{u}\cdot\vec{v} = u_x v_x + u_y v_y + u_z v_z = \|\vec{u}\|\,\|\vec{v}\|\cos\theta
    \]
    \textbf{Orthogonalité :} $\vec{u} \perp \vec{v} \iff \vec{u}\cdot\vec{v} = 0$.
    \item \textbf{Produit vectoriel $\vec{n} = \vec{u} \wedge \vec{v}$ ($\perp \vec{u}$ et $\perp \vec{v}$) :}
    \[
        \vec{u}\wedge\vec{v} = \bigl(u_y v_z - u_z v_y,\;\; u_z v_x - u_x v_z,\;\; u_x v_y - u_y v_x\bigr)
    \]
    \item \textbf{Coplanarité de 3 vecteurs $\vec{u}, \vec{v}, \vec{w}$ :}
    \[
        \det(\vec{u},\vec{v},\vec{w}) = \begin{vmatrix} u_x & v_x & w_x \\ u_y & v_y & w_y \\ u_z & v_z & w_z \end{vmatrix} = 0
    \]
\end{itemize}

\vspace{4pt}
\noindent{\large\bfseries\color{uclblue} 2. Droites et Plans dans l'Espace}
\begin{itemize}[leftmargin=1.2em, itemsep=2pt, topsep=2pt]
    \item \textbf{Droite $d$ passant par $A$ de vecteur directeur $\vec{u} = (a,b,c)$ :}
    \[
        d \equiv \begin{cases} x = x_A + t\,a \\ y = y_A + t\,b \\ z = z_A + t\,c \end{cases} \quad (t \in \R)
    \]
    Distance parcourue entre $A(0)$ et $M(t)$ : $|AM| = |t|\cdot\|\vec{u}\|$.
    \item \textbf{Plan $\Pi$ de vecteur normal $\vec{n} = (a,b,c)$ :}
    \[
        \Pi \equiv a(x-x_A) + b(y-y_A) + c(z-z_A) = 0
    \]
    soit $\mathbf{ax + by + cz + d = 0}$ avec $\vec{n} = (a,b,c) \perp \Pi$.
    \item \textbf{Plan défini par un point $A$ et deux vecteurs directeurs $\vec{u}, \vec{v}$ :} Prendre $\vec{n} = \vec{u} \wedge \vec{v}$ ou résoudre $\det(\overrightarrow{AM},\vec{u},\vec{v}) = 0$.
    \item \textbf{Plan médiateur d'un segment $[AB]$} (lieu des points équidistants de $A$ et $B$) : Plan passant par le milieu $M = \frac{A+B}{2}$ et de vecteur normal $\vec{n} = \overrightarrow{AB}$.
\end{itemize}

\vspace{4pt}
\noindent{\large\bfseries\color{uclblue} 3. Positions Relatives Droites / Plans}
Soit $d$ de vecteur directeur $\vec{u}$ et $\Pi_1, \Pi_2$ de normales $\vec{n}_1, \vec{n}_2$ :
\begin{itemize}[leftmargin=1.2em, itemsep=2pt, topsep=2pt]
    \item \textbf{Droite parallèle au plan ($d \parallel \Pi_1$) :} $\vec{u} \cdot \vec{n}_1 = 0$
    \item \textbf{Droite perpendiculaire au plan ($d \perp \Pi_1$) :} $\vec{u} \parallel \vec{n}_1$
    \item \textbf{Plans perpendiculaires ($\Pi_1 \perp \Pi_2$) :} $\vec{n}_1 \cdot \vec{n}_2 = 0$
    \item \textbf{Plans parallèles ($\Pi_1 \parallel \Pi_2$) :} $\vec{n}_1 \parallel \vec{n}_2$
\end{itemize}

\columnbreak

\noindent{\large\bfseries\color{uclblue} 4. Intersections, Projections et Distances}
\begin{itemize}[leftmargin=1.2em, itemsep=2pt, topsep=2pt]
    \item \textbf{Point de percée $B = d \cap \Pi$ :} Remplacer $x(t), y(t), z(t)$ de la droite $d$ dans l'équation cartésienne $ax+by+cz+d=0$ du plan $\Pi$ pour trouver $t$.
    \item \textbf{Projeté orthogonal $P$ d'un point $C$ sur une droite $d \equiv A + t\vec{u}$ :} Écrire $P(t) = A + t\vec{u}$ et résoudre l'équation d'orthogonalité :
    \[
        \overrightarrow{CP(t)} \cdot \vec{u} = 0
    \]
    \item \textbf{Projeté orthogonal $H$ d'un point $A$ sur un plan $\Pi$ de normale $\vec{n}$ :} $H$ est le point de percée dans $\Pi$ de la droite passant par $A$ et de vecteur directeur $\vec{n}$.
    \item \textbf{Droite $d_2$ parallèle à $d_1 \equiv A + t\vec{u}$ à distance $D$ dans la direction $\vec{v} \perp \vec{u}$ :}
    \[
        B = A + D\,\frac{\vec{v}}{\|\vec{v}\|} \quad \text{et} \quad d_2 \equiv B + t\,\vec{u}
    \]
    \item \textbf{Distance d'un point $A(x_A, y_A, z_A)$ au plan $\Pi \equiv ax + by + cz + d = 0$ :}
    \[
        d(A, \Pi) = \frac{|ax_A + by_A + cz_A + d|}{\sqrt{a^2 + b^2 + c^2}}
    \]
\end{itemize}

\vspace{4pt}
\noindent{\large\bfseries\color{uclblue} 5. Réflexion Optique sur un Miroir Plan}
Un rayon lumineux de vecteur directeur incident $\vec{v}_{\text{inc}}$ frappant un miroir plan $\Pi$ de vecteur normal $\vec{n}$ en un point $B$ fait un angle d'incidence $\theta$ avec la normale ($\cos\theta = \frac{|\vec{v}_{\text{inc}}\cdot\vec{n}|}{\|\vec{v}_{\text{inc}}\|\|\vec{n}\|}$) et repart selon le vecteur directeur réfléchi :
\[
    \vec{v}_{\text{réf}} = \vec{v}_{\text{inc}} - 2\,\frac{\vec{v}_{\text{inc}} \cdot \vec{n}}{\|\vec{n}\|^2}\,\vec{n}
\]

\vspace{4pt}
\noindent{\large\bfseries\color{uclblue} 6. Sphère et Section par un Plan}
\begin{itemize}[leftmargin=1.2em, itemsep=2pt, topsep=2pt]
    \item \textbf{Sphère $\mathcal{S}$ de centre $A(x_A, y_A, z_A)$ et de rayon $R$ :}
    \[
        (x - x_A)^2 + (y - y_A)^2 + (z - z_A)^2 = R^2
    \]
    \item \textbf{Intersection de $\mathcal{S}(A, R)$ avec un plan $\Pi$ :} Si $d = d(A, \Pi) < R$, l'intersection est un \textbf{cercle} (ou un disque pour la boule) centré au projeté orthogonal $H$ de $A$ sur $\Pi$, de rayon :
    \[
        r = \sqrt{R^2 - d^2} \quad \implies \quad \text{Aire} = \pi(R^2 - d^2)
    \]
    \item \textbf{Plan tangent à $\mathcal{S}(A,R)$ en un point $T \in \mathcal{S}$ :} C'est le plan passant par $T$ et de vecteur normal $\vec{n} = \overrightarrow{AT}$.
\end{itemize}

\end{multicols}
}

\restoregeometry
